\chapter{Introducción}

\section{Motivación}

Un alumno con discapacidad auditiva que asiste a una clase magistral se enfrenta a un
problema que el resto no tiene: la información llega por un canal que no puede usar. Las
soluciones habituales pasan por un intérprete de lengua de signos o por un profesional que
transcriba en directo, y ambas comparten la misma limitación: dependen de que haya una
persona disponible, con formación específica, durante toda la sesión. En la práctica eso
significa que muchas clases se quedan sin cobertura.

Hay un segundo grupo, mucho más numeroso, para el que el canal tampoco funciona del todo:
el alumnado que todavía no domina la lengua en la que se imparte la clase. En el sistema
educativo español no universitario había en el curso 2023-2024 más de un millón de alumnos
de nacionalidad extranjera, el 12,2\% del total \cite{mefpd2024datos}; una parte
considerable procede de países donde el castellano no es lengua vehicular. Su dificultad no
es sensorial sino de velocidad: entienden el vocabulario cuando lo leen y lo pierden cuando
lo oyen a la velocidad del habla espontánea, con la fonética de un hablante nativo y sin
posibilidad de volver atrás.

Los dos casos difieren en casi todo menos en lo que necesitan, que es exactamente lo mismo:
que lo dicho aparezca escrito, en el momento, sin pedírselo a nadie. Y la evidencia de que
eso ayuda al segundo grupo no es una suposición razonable: está medida, y se recoge en el
capítulo~\ref{cap:estado}. Conviene decir desde ya que este trabajo \emph{no} ha evaluado
el sistema con alumnado de ninguno de los dos colectivos; la segunda población se menciona
porque amplía el alcance de lo construido, no porque el trabajo aporte evidencia sobre
ella.

El reconocimiento automático del habla ha alcanzado en los últimos años una calidad que
invita a plantear la alternativa. Los modelos de propósito general funcionan
razonablemente bien en español sin necesidad de entrenarlos, y el equipamiento necesario
para ejecutarlos ha dejado de ser exclusivo de centros de cálculo. La pregunta ya no es si
la transcripción automática es viable, sino si lo es \emph{con la calidad suficiente} para
que un alumno pueda seguir una clase sin perder contenido.

Esa distinción es la que motiva este trabajo. Un sistema que acierta el noventa por ciento
de las palabras suena aceptable expresado así; pero si el diez por ciento restante incluye
los términos técnicos de la asignatura, los nombres propios o las negaciones, el alumno no
recibe una versión ligeramente degradada de la clase: recibe una versión distinta. La
métrica habitual del campo, la tasa de error de palabra, no distingue entre perder una
preposición y perder un \emph{no}.

\section{Planteamiento del trabajo}

El trabajo tiene dos núcleos que se abordan por separado y se apoyan mutuamente.

El primero es una comparativa empírica: se evalúa si las técnicas de adaptación al dominio
más habituales mejoran el reconocimiento en español sobre material que se parece al de un
aula. Se estudian tres: aportar al modelo un glosario de la asignatura como contexto,
corregir su salida con un modelo de lenguaje, y ajustar sus pesos mediante LoRA. Cada una
se mide con el mismo diseño, sobre el mismo corpus y con el mismo criterio estadístico, de
modo que las cifras sean comparables entre sí.

El segundo es una aplicación de subtitulado en vivo, construida para desplegarse en la red
local de un centro: el docente emite desde su equipo y el alumnado lee desde el suyo, sin
que ningún dispositivo del aula necesite capacidad de cálculo. La aplicación no es una
demostración del primer núcleo; es un banco de pruebas que plantea preguntas propias, como
cada cuánto trocear el audio o dónde cortarlo, que se responden con la misma metodología
experimental.

La relación entre ambos resultó menos unidireccional de lo previsto. La investigación
proporcionó a la aplicación criterios de diseño medidos en lugar de intuidos; pero la
aplicación detectó a su vez un problema que ningún experimento de laboratorio podía ver,
porque solo aparece cuando hay un micrófono y un navegador reales de por medio.

\section{Aportaciones}

Conviene adelantar a qué se ha llegado, porque condiciona cómo debe leerse el resto. Cinco
resultados resumen el trabajo.

\begin{enumerate}
  \item \textbf{De las tres técnicas de adaptación, solo una mejora.} El ajuste fino con
        LoRA reduce el error de palabra 1,23 puntos con dos criterios estadísticos de
        acuerdo. El prompting contextual no produce efecto detectable, y la post-corrección
        con un modelo de lenguaje \emph{degrada} de forma significativa y replicada sobre
        dos corpus.
  \item \textbf{El margen no está donde se suponía.} Segmentar el audio por silencios en
        lugar de por reloj reduce el error 7,4 puntos a igual latencia media, con los dos
        mismos criterios estadísticos de acuerdo: seis veces más que la mejor técnica de
        adaptación, y a coste de cómputo nulo. La premisa de partida del trabajo queda
        invertida por sus propios datos.
  \item \textbf{El WER agregado oculta lo que importa en accesibilidad.} Con un 18,61\% de
        error, cifra que suena tolerable, el sistema pierde o altera el 28\% de las
        negaciones e introduce trece que nadie pronunció. Se proponen dos métricas
        complementarias que exponen esos fallos.
  \item \textbf{El modo de fallo pesa más que la tasa de error al elegir modelo.} El
        criterio descartó primero a la variante más rápida de la familia elegida, y después
        a un candidato de arquitectura distinta que ganaba en las métricas pero traducía al
        inglés ante habla española difícil.
  \item \textbf{Varias mediciones aparentemente correctas eran falsas.} Se documentan ocho
        incidentes en los que el procedimiento producía cifras verosímiles y erróneas sin
        emitir ningún error, y las comprobaciones que se derivaron de ellos.
\end{enumerate}

En el plano aplicado, el trabajo entrega además un sistema de subtitulado desplegable en la
red de un centro, con el reconocimiento centralizado en el servidor, y en el que la técnica
ganadora se incorporó sin modificar una sola línea de la aplicación.

\section{Estructura del trabajo}

El capítulo~\ref{cap:estado} sitúa el trabajo respecto a la literatura y a la práctica del
sector, y documenta la búsqueda de corpus, que resultó ser el condicionante principal de
todo el diseño experimental. El capítulo~\ref{cap:requisitos} deriva los requisitos, tanto
funcionales de la aplicación como metodológicos de la comparativa. El
capítulo~\ref{cap:objetivos} enuncia los objetivos y los criterios con los que se
consideran cumplidos.

El capítulo~\ref{cap:desarrollo} contiene el grueso del trabajo: la metodología
experimental, los diez experimentos realizados con sus resultados, y el desarrollo de la
aplicación. El capítulo~\ref{cap:conclusiones} recoge las conclusiones, la verificación de
los objetivos, las limitaciones reconocidas y las líneas que quedan abiertas. El
apéndice~\ref{ap:reproducibilidad} documenta cómo regenerar cualquier cifra de la memoria,
el apéndice~\ref{ap:riesgos} el registro de riesgos que guió la planificación, y el
apéndice~\ref{ap:despliegue} los pasos concretos para poner el sistema en marcha en un
centro.

Conviene anticipar una particularidad de este trabajo: buena parte de sus aportaciones no
proviene de que las técnicas estudiadas funcionaran, sino de descubrir bajo qué condiciones
las mediciones que las evalúan son válidas. Varias de las conclusiones surgieron de
detectar que un resultado aparentemente sólido era un artefacto del procedimiento de
medida. Esos hallazgos se documentan con el mismo detalle que los resultados positivos,
porque condicionan cómo deben interpretarse todos los demás.
